\SourceLeaf{337}

The signs of
\[ \sin(\mathfrak{B}-B), \qquad \sin(\mathfrak{C}-C), \qquad \sin(\mathfrak{B}+\mathfrak{C}-B-C) \]
in the eight possible combinations will be as follows.

\begin{center}
\begin{tabular}{cc|cc}
I 1 & $+\;+\;+$ & II 1 & $-\;-\;-$ \\
I 2 & $-\;+\;-$ & II 2 & $+\;-\;+$ \\
I 3 & $+\;-\;-$ & II 3 & $-\;+\;+$ \\
I 4 & $-\;-\;+$ & II 4 & $+\;+\;-$ \\
\end{tabular}
\end{center}

Thus it is necessary and sufficient for the physical possibility of the problem that the three sines have the same sign.

If the first sine has a sign different from the other two, then (case 3) the construction gives a point lying in the unsuitable segment of the circle $\gamma$.

If the second sine has a sign different from the other two, then [case 2] the point found by the construction lies in the unsuitable segment of the circle $\beta$.

\Needspace{6\baselineskip}If instead the third sine differs in sign from the other two, then [case 4] the point found by the construction lies in unsuitable segments of both circles $\beta$ and $\gamma$.

\hfill 1852 July 29.

\medskip

\begin{center}
\textbf{REMARKS.}
\end{center}

Most of the preceding notes are in a notebook, with only [7] on a separate sheet. As for when they were written, note [1] dates from 1832, whereas [5] and [6] are to be dated to 1849, and [7] to 1852.

From December 1839 until Easter 1840 \textsc{Gauss} gave a course of lectures on the \emph{Theory of Imaginary Quantities}, which E. \textsc{Heine} attended and wrote up. The following discussions of the quadrilateral are taken from his write-up and may help explain the notes gathered under [III]. \hfill \textsc{Stäckel}.
